\section{A universal relative golden-head bound and hereditary full C}
Put q=phi-1, psi=-q, so q²+q=1. For EVERY arbitrary ordered
nonempty nonbinary reference V, let a=sum V, R its infinite-BINARY
reference, and H_j=R_j-phi R_(j-1), H_0=1. We prove the UNIFORM
relative head bound
\[
                   H_j\ge q^6 H_a>0\qquad(2<=j<=a).     \tag{1}
\]
The number of factors, all capacities and a are unbounded.
There is no factor-count or weight-dependent constant in (1).

\paragraph{Normalize at the true terminal gap.}
The exact homogeneous tail is H_(a+m)=psi^m H_a. At an append t>=3,
\[
 H'_(a+t)=H_a\sum_{m=0}^{t-1}\psi^m
        =q(1-\psi^t)H_a.
\]
Thus H_a=rho=prod_(t in V) q(1-psi^t), with0<rho<1;
every append contracts rho since each factor<=2q²<1.
The first-capacity head has H_j=q+psi^(j+1)>=q², and (1) holds.

For further appends use the exact full interval
H'_n=sum_(j=max(0,n-t)..n-1)H_j. Every alternating tail group
starting at a has sum at least q²rho. Set c=q^6.
If its left endpoint is at least2 and below a, it contains either
a nonempty old positive head of total at least c*rho, or this
head followed by a nonnegative tail group. If it starts at a,
the tail group alone gives q²rho>=c*rho_new.
The normalized endpoint contracts, so c*rho>=c*rho_new.

A window starting at0 retains H0+H1=1-q>=c. A window starting
at1 and with a>=4 retains H1+H2+H3=3q²-q>=c before any tail.
The remaining a=3 case retains exactly
\[
                  -q+q²+q²H_3=-q+q²+2q^4=q^6=c.
\]
Since rho_new<1 this also pays c*rho_new. These cases cover
every legal new current index2..a+t and prove (1) by complete
word induction. The entire quantitative interval-window induction,
not only its scalar end payment, is included in the new Lean file.

\paragraph{Full reference C survives every capacity doubling of nonbinary weight.}
Suppose the old SORTED reference V already has full C. Append
any t>=2a. Conditional old-C transfer gives new C through t+4,
by finite-prefix realization; the universal reference theorem
gives it from a+t-1. For a remaining n with t+4<n<a+t, put
s=n-t. Then2<=s-1<=a, and n-1>=a+7. Consequently (1) gives
\[
 H_(s-1)\ge q^6\rho,
 \qquad H_(n-1)\le q^7\rho<q^6\rho.
\]
The EXACT new derivative golden gap is H_(s-1)-H_(n-1)>0.
The universal reference cumulative golden margin pays EVERY new
C read. This proves full C for the extended reference at all ranks.

More generally (1) is an unconditional earlier reference resource:
for any last append t, every current with n>=max(t+3,a+8)
and n<=a+t+1 has a positive golden-D gap and hence all C shifts.
This conclusion does not require old full C. Empty read ranges
and the known universal reference tail are understood separately.

\paragraph{Full ACTUAL C survives the same append, uniformly in binary count.}
Assume that V is a fixed nonempty sorted tail for which ALL
actual1^alpha2^uV have full C, for EVERY alpha>=1 and u>=0.
Finite-prefix realization also gives old reference full C, so
the preceding argument gives the extended reference full C.
Let h=a-r, r=len V, append t>=2a, and fix arbitrary alpha,u.
The conditional C prefix/mode theorem and universal actual tail
leave only t+4<n<a+t, n<=ell, where ell is the OLD first D mode.
The general center theorem gives
\[
                         2ell<=3u+a+r+1.                 \tag{2}
\]
If u<h, (2) gives ell<=t+4. If u+r<=5, the same conclusion
follows from a>=3,r>=1 and t>=2a. These cases are closed.

We may therefore assume u>=h and u+r>=6. Put L=u+r+1>=a+1.
The recent surplus ordinary theorem gives F LC from L; its
entire prefix through L is protected relative to R. The NEW
prefix through L+1 is protected relative to the extended reference.
Thus n<=L+1 is paid by its complete C theorem.

For n>=L+2 the exact new strictly growing positive derivative pair is
A=F_(n-2)-F_(n-t-2)>0, B=F_(n-1)-F_(n-t-1)>A.
Its outgoing current index n-t-1 lies in2..a-2<=L and has its
protected ordinary-golden LOWER bound. The incoming index n-1
is at least max(a,L+1), and has its ACTUAL ordinary-golden CEILING.
Therefore mu=B/A<F_(n-1)/F_(n-2)<phi. All old G gaps through L
have their protected q² margin; ordinary F LC from L integrates
this margin through n-2. The exact affine G window, including its
rank-zero atom1, pays every new G read at mu because t>=3 and
2q²>=q. This excludes every remaining comparison and proves
complete actual C uniformly in alpha,u. No normalized actual
relative-LC or density-record hypothesis is used.

\paragraph{Arbitrarily many nonbinary factors, all binary counts.}
Start with ANY tail V for which uniform full C is already known.
Append any finite sequence of capacities satisfying
\[
                  t_j\ge2\sum_{i<j}t_i                 \tag{3}
\]
at EACH new step. The preceding theorem and induction prove
complete actual C for every resulting tail and EVERY binary count
and initial-one length. The factor count and all capacities are
unbounded. In particular one may start with two completely arbitrary
3<=L<=T, using the new full two-nonbinary theorem, or ANY old
weight-at-most32 tail, using the prior complete weight theorem.
The condition concerns only nonbinary weight, with NO bound
depending on binary count or initial padding.

This full-C theorem supplies ordinary-product coefficient-tilt UI
against any opposite arm and the registered every-depth BOTH-rising
original C-to-S nonrecovery, with current/physical guards and OLD
weighted-A2 head premise retained. It does not prove unrestricted
A2 or original core coefficient UI. General nondominant three-or-
more-nonbinary full C remains OPEN, as does the original conjecture.

Eight scoped Lean reports prove the all-q atom identities, full
quantitative interval-window theorem, ALL dominant-current reads,
ALL deficit/small-count mode exclusion, surplus read guards,
strict tail gap and complete natural-index predicate assembly.
Reference/actual word, CI/norm, Fibonacci-tail and source/coupled
identifications and the full tail-word induction are written using
registered inputs. No end-to-end Lean composition datatype theorem
is asserted. All finite checks are QA only; no new finite premise.
